\documentclass[11pt]{article}
\usepackage[T1]{fontenc}
\usepackage{lmodern}
\usepackage{amsmath,amssymb}
\usepackage{booktabs}
\usepackage{array}
\usepackage[a4paper,margin=2.5cm]{geometry}
\usepackage[hidelinks]{hyperref}
\newcommand{\doi}[1]{\href{https://doi.org/#1}{\texttt{#1}}}
\setlength{\parskip}{4pt}
\title{\vspace{-1.6cm}\Large Guided Lamb Waves in a One-Dimensional Hexagonal Piezoelectric Quasicrystal Plate\\[2pt]
\large with Phason Relaxation, a Size Effect and a Single-Lag Heat Equation\\[8pt]
\large Research plan, first half}
\author{Dr.\ Vipin \\ Gurugram University}
\date{30 October 2026}
\begin{document}
\maketitle
\vspace{-6mm}

\section{Title, abstract, keywords}

\paragraph{Working title} Lamb waves in a one-dimensional hexagonal piezoelectric quasicrystal plate:
phason relaxation, nonlocal strain gradient elasticity and thermal coupling in one dispersion problem.

\paragraph{Abstract} We study plane-strain Lamb waves in a plate of thickness $2h$ cut from a
one-dimensional hexagonal piezoelectric quasicrystal, with free faces at $z=\pm h$ and the quasiperiodic
direction along the thickness. The model carries five fields: the two phonon displacements, the phason
displacement, the electric potential and the temperature rise. The phason field obeys a single equation that
contains both an inertia and a friction term, so the wave-like and the diffusive descriptions of phason
motion are the two limits of one parameter rather than two competing models. The phonon stress is written in
the nonlocal strain gradient form, with the two size lengths entering only as length over thickness, and
every piezoelectric, thermal and phonon-phason term is kept local. The heat equation is of the modified
gradient type in its base form, with one phase lag, and it carries the work of the phason stress as well as
the work of the phonon stress. The paper asks three things. Which of the two phason limits survives in a
plate of finite thickness, and where the crossover sits. How much attenuation the phason friction adds to
each Lamb branch. How the thermal relaxation competes with the phason relaxation, since both are times and
both cut off the high frequency behaviour. The plan gives the model, the sequence of comparisons to be made
before any claim, and the figures the paper will carry.

\paragraph{Keywords} quasicrystal; phason; Lamb waves; piezoelectric plate; nonlocal strain gradient
elasticity; attenuation.

\section{Introduction}

A quasicrystal is a solid with long-range order but no periodicity. Its diffraction pattern has sharp spots,
which says that the arrangement repeats in a certain sense, but the spots sit at positions that no periodic
lattice can produce. The standard way to describe such a material in continuum mechanics is to give it a
second displacement field. The phonon displacement is the ordinary displacement of the material points. The
phason displacement is not a displacement in space at all; it measures how the tiles of the aperiodic tiling
rearrange locally, sliding over one another without changing their shape. A gradient of the phason
displacement produces a phason strain, and a phason strain costs energy. This framework goes back to the
mid-eighties and the hydrodynamics of the icosahedral phase by Lubensky, Ramaswamy and Toner
\cite{Lubensky1985}, and to the elasticity of pentagonal and icosahedral quasicrystals by Levine, Lubensky,
Ostlund, Ramaswamy, Steinhardt and Toner \cite{Levine1985}. Both of those papers are about the icosahedral
case, and the moduli in them should not be transferred to a hexagonal quasicrystal; we cite them for the
concept and for the relaxation mechanism.

One-dimensional hexagonal quasicrystals are the simplest case and the only one with a mature plate
literature. The quasiperiodicity runs along a single axis, which in our problem is the thickness. Guided
waves in such plates have been computed by Zhang, Yu, Zhang and Elmaimouni \cite{Zhang2021}, and the same
group has since added nonlocal effects with an integral kernel \cite{Wang2025}, couple stress \cite{Wang2022}
and layered construction \cite{Wang2025b}. The message from that body of work is that the phonon-phason
coupling shifts the branches in a way that is measurable, and that the nonlocal parameter and the thickness
compete in a way that can be mapped.

Two things are missing from it, and both are about time. The first is the dynamics of the phason itself.
Most plate papers treat the phason as if it were in equilibrium, because the phason relaxes. The
distinction matters when the relaxation time is comparable with the wave period, which is the case the
phason wave parameter study of Chellappan, Gopalakrishnan and Mani \cite{Chellappan2015} puts on a formal
footing for bulk quasicrystals. That is a
choice, not a law. Agiasofitou and Lazar wrote the phason momentum balance with both inertia and a friction
tensor and showed that the diffusion limit and the wave limit are the two ends of one telegraph equation
\cite{Agiasofitou2014}. Baggioli and Landry put the same crossover in the compact form in which the phason
frequency obeys a damped wave relation, diffusive at long wavelength and propagating at short
\cite{Baggioli2020}. Neither paper is about plates. A plate of finite-thickness phason modes is exactly the
regime in which the crossover should be visible, because the thickness sets the smallest wavenumber the
plate can hold.

The second missing thing is heat. A quasicrystal plate at a different temperature from its surroundings
carries a temperature field, and both the phonon and the phason stresses do work against it. Whether the
heat equation of a quasicrystal plate should carry a relaxation time is not settled in the literature we
have read, and the closest paper, by Zhang, Song, Li, Han, Yu and Yang on generalized thermoelastic Lamb
waves in hexagonal quasicrystal nanoplates \cite{Zhang2026}, must be read before we write a single sentence
about novelty. What we can say now is what we will do differently: we put both phason terms in one equation,
we carry the phason work in the heat equation, and we report attenuation rather than only dispersion.

To our knowledge, after searching Crossref and the open web in October 2026 for Lamb waves in
one-dimensional hexagonal piezoelectric quasicrystal plates, no paper carries a telegraph phason equation
together with a relaxation-time heat equation in one dispersion problem. That sentence is the whole novelty
claim, and it is narrower than it looks. If the 2026 paper already carries the same heat operator, then the
heat part of our contribution disappears and what remains is the phason inertia and friction pair and the
attenuation curves. Neha should read that paper in week one and write the answer in her notebook before she
writes anything else.

\section{Problem statement and assumptions}

The plate has thickness $2h$ and free faces at $z=\pm h$. The quasiperiodic direction is $z$. Motion is
plane strain in the $x$--$z$ plane. Fields vary as $\exp[i(kx-\omega t)]$ with $k=k_r+ik_i$, and a positive
$k_i$ means decay along $+x$. Frequency is reported through $f\,(2h)$. The five unknowns are the phonon
displacements $u_x$ and $u_z$, the phason displacement $w_z$, the electric potential $\varphi$ and the
temperature rise $\theta$.

\begin{enumerate}
\item Linear phonon and phason elasticity, with constant moduli at the reference temperature.
\item Small phason strain, so that the phason stress is linear in the phason strain and in the phonon
strain.
\item A single phason inertia $\rho_p$ and a single phason friction $\gamma$, taken isotropic in the plane
of the plate. Anisotropic friction is a labelled sensitivity case, not the main model.
\item Piezoelectric coupling as in a transversely isotropic hexagonal crystal with the axis along $z$.
\item Small temperature rise, $|\theta|/T_0\ll1$, so the heat equation can be linearised about the
reference state.
\item No free carriers; the electric problem is a dielectric one, with a possible surface charge condition
at the faces which the scholar will choose.
\end{enumerate}

Two of these deserve a sentence. The phason friction is the one parameter in the model that no textbook
gives for a plate, and it must be swept rather than fixed, with the sweep reported. And the assumption of a
single inertia and a single friction is a simplification of the tensor form in Agiasofitou and Lazar
\cite{Agiasofitou2014}; we take it because the plate problem already has enough parameters and because the
anisotropy of friction is not something a wave test can resolve.

\section{The model}

\subsection{Constitutive relations}

With $\varepsilon_{ij}=\tfrac12(u_{i,j}+u_{j,i})$, $w_{ij}=w_{i,j}$ and $E_i=-\varphi_{,i}$, the phonon
stress carries the size effect on its elastic part only,
\begin{equation}
\bigl(1-l_{\rm nl}^2\nabla^2\bigr)t_{ij}=\sigma_{ij},\qquad
\sigma_{ij}=C_{ijkl}\bigl(1-l_g^2\nabla^2\bigr)\varepsilon_{kl}+R_{ijkl}w_{kl}
-e_{kij}E_k-\beta_{ij}\theta ,
\end{equation}
the phason stress is local and carries the phonon strain and the phason strain,
\begin{equation}
H_{ij}=R_{klij}\varepsilon_{kl}+K_{ijkl}w_{kl}-\beta^{w}_{ij}\theta ,
\end{equation}
and the electric displacement is local,
\begin{equation}
D_i=e_{ikl}\varepsilon_{kl}+\epsilon_{ik}E_k+p_i\theta .
\end{equation}
The two operators act on the phonon elastic strain energy and nowhere else. The phonon-phason coupling, the
piezoelectric coupling, the pyroelectric coupling and both thermal stresses are local. The reason is the
same as in the companion paper: keeping $D_i$ second order means the electric problem needs no extra
condition at the faces, and keeping the size effect mechanical means the extra conditions are mechanical
too.

The thermo-phason term $\beta^{w}_{ij}\theta$ is an extension, not a published relation. Duhamel's analogy
extends to the phason stress in the obvious way, but we have found no paper that writes it, and the paper
must say so in the sentence where it appears. If Neha cannot find a source, the term stays with an explicit
statement of its status.

\subsection{Field equations}

Momentum, in the form in which the nonlocal operator acts on inertia,
\begin{equation}
\bigl(1-l_{\rm nl}^2\nabla^2\bigr)\rho\,\ddot u_i=\sigma_{ij,j},
\end{equation}
the phason balance with inertia and friction,
\begin{equation}
H_{zj,j}=\rho_p\ddot w_z+\gamma\dot w_z ,
\end{equation}
Gauss law in a dielectric,
\begin{equation}
D_{i,i}=0 ,
\end{equation}
and the heat equation,
\begin{equation}
\Bigl(1+\tau_q\frac{\partial}{\partial t}\Bigr)
\Bigl[\rho c_p\dot\theta+T_0\beta_{ij}\dot\varepsilon_{ij}
+T_0\beta^{w}_{ij}\dot w_{ij}+T_0p_i\dot E_i\Bigr]
=k_{\rm th}\nabla^2\theta .
\end{equation}
Five equations for five unknowns. Equation four is the one this paper is about. Setting $\gamma=0$ leaves a
wave equation for the phason with a finite speed; setting $\rho_p=0$ leaves a diffusion equation; keeping
both leaves a telegraph equation whose dispersion runs from one limit to the other as the wavenumber grows.
In a plate the allowed wavenumbers are set by the thickness, so the crossover can be moved through the
measured range by changing the thickness. That is the experiment the paper proposes, in principle.

\subsection{Notation}

\begin{center}
\small
\begin{tabular}{p{2.4cm}p{6.7cm}p{3.4cm}}
\toprule
symbol & meaning & unit\\
\midrule
$u_x,u_z$ & phonon displacements & m\\
$w_z$ & phason displacement along the quasiperiodic axis & m\\
$\varphi$ & electric potential & V\\
$\theta$ & temperature rise over $T_0$ & K\\
$k=k_r+ik_i$ & wavenumber; $k_i>0$ is decay along $+x$ & m$^{-1}$\\
$2h$ & plate thickness; faces at $z=\pm h$ & m\\
$C_{ijkl}$ & phonon elastic moduli & Pa\\
$K_{ijkl}$ & phason elastic moduli & Pa\\
$R_{ijkl}$ & phonon-phason coupling moduli & Pa\\
$e_{kij}$, $\epsilon_{ik}$, $p_i$ & piezoelectric, dielectric, pyroelectric constants & C\,m$^{-2}$, F\,m$^{-1}$, C\,m$^{-2}$K$^{-1}$\\
$\rho$, $\rho_p$ & mass density, phason mass density & kg\,m$^{-3}$\\
$\gamma$ & phason friction coefficient & kg\,m$^{-3}$s$^{-1}$\\
$l_{\rm nl}$, $l_g$ & nonlocal and gradient lengths, used as $l/h$ & m\\
$\beta_{ij}$, $\beta^{w}_{ij}$ & thermal stress, thermo-phason coupling & Pa\,K$^{-1}$, Pa\,K$^{-1}$\\
$\rho c_p$, $k_{\rm th}$, $\kappa$ & heat capacity per volume, conductivity, diffusivity & J\,m$^{-3}$K$^{-1}$, W\,m$^{-1}$K$^{-1}$, m$^2$s$^{-1}$\\
$\tau_q$, $\tau_T$ & phase lag, $\tau_T=(2h)^2/(\pi^2\kappa)$ & s\\
$\omega$, $f$ & angular frequency, frequency & rad\,s$^{-1}$, Hz\\
\bottomrule
\end{tabular}
\end{center}

\subsection{Order of the system, and where this chapter stops}

The nonlocal and gradient operators raise the order of the problem in $z$, and the phason equation adds a
second-order field with no classical counterpart, so the free faces need more conditions than the classical
traction and electrical pair, and some of those extra conditions will be phason conditions. Which ones to
use, and whether they should be phrased in terms of the phason stress or the phason displacement, is your
problem to settle; this chapter stops before the faces.

\section{Work plan}

Six steps, each ending with a written comparison.

\paragraph{Step 1. The elastic quasicrystal plate, both limits of the phason} With $e=0$, $\theta=0$ and both
size lengths zero, the model must reproduce the published dispersion of a one-dimensional hexagonal
quasicrystal plate \cite{Zhang2021} in shape and ordering, with the phonon-phason coupling switched off and
then on. With $\gamma=0$ the phason branch must be a propagating branch at a speed set by the phason moduli;
with $\rho_p=0$ it must be a diffusive branch that does not propagate. Both are computed before any coupled
case, and the branch count is written down.

\paragraph{Step 2. Piezoelectric coupling} Switch on $e_{kij}$ and compare the branch splitting with the
piezoelectric quasicrystal plate results of the same group \cite{Wang2025,Wang2022}. The comparison is of
ordering and of the sign of the shift, since the material set is ours.

\paragraph{Step 3. Heat in the two limits} With the phason friction set aside for the moment, switch on $\theta$ and
compare the isothermal and adiabatic limits. The coupled answer must lie between them, to a width of order
the thermoelastic coupling. Report the width, and report separately how much of it comes from the phason
work term, which is the one we have no source for.

\paragraph{Step 4. Two times in competition} With both the phason friction and the thermal lag active, the
high-frequency behaviour is cut off twice. The heat side cuts off at $1/|\tau_T-\tau_q|$; the phason side
cuts off when the phason frequency falls below the friction rate. The paper must show which cutoff comes
first for the thicknesses of interest. This is the figure that decides whether the thermal field is needed at
all, and it is drawn before the parameter sweep, not after.

\paragraph{Step 5. Switching-off tests} Set $R_{ijkl}=0$ and the phonon-phason coupling must vanish, with the
phonon and phason problems decoupled to within $10^{-6}$ relative on $k_r$ and $k_i$ for every sampled
mode. Set $\gamma$ to zero and the result must match an independently assembled lossless phason model to
within the same tolerance; set $\rho_p$ to zero and it must match the diffusion model. Set the thermal
coupling to zero and the attenuation must fall to the mechanical value. Every residual goes in an appendix
table. A tolerance of $10^{-4}$ is a warning, not a pass.

\paragraph{Step 6. Size effect and the final sweep} Sweep $l_{\rm nl}/h$ and $l_g/h$ as in the companion paper,
then $\gamma$ over the range that keeps the phason branch on both sides of the crossover, and produce the
attenuation families for three thicknesses. Report the phason contribution as a difference of attenuations.

\section{Figures, risks and timeline}

Four figures. First, the dispersion of the two phason limits on the same axes, with the propagating branch
and the diffusive branch labelled, which is the clearest statement of what the paper is doing. Second, the
attenuation of the two lowest branches against frequency times thickness, with the isothermal and adiabatic
brackets as bands. Third, the competition of the two cutoff times as a map of thickness against friction
rate, with the boundary where the thermal cutoff comes first drawn on it. Fourth, the sensitivity of the
lowest antisymmetric branch to the nonlocal length, since that is the parameter a referee will press on.

The first risk is novelty and it is serious. The 2026 paper on generalized thermoelastic Lamb waves in
hexagonal quasicrystal nanoplates \cite{Zhang2026} may already contain the heat operator we propose to use.
Nothing in this plan can be fixed by writing more; the paper has to be read in week one, and if it closes the
thermal route then the paper narrows to the phason telegraph equation and the attenuation curves. The second
risk is that the phason friction cannot be measured in a plate, so the sweep may be dismissed as arbitrary;
the answer is to report the friction as a range tied to a measured phason relaxation time, where one exists
for a real material. The third risk is that the two-size-length description of a quasicrystal is not on firm
ground, since the nonlocal kernel used in the closest papers is integral and ours is differential; the
comparison in Step 1 is designed to expose any inconsistency. The fourth risk is the missing source for the
thermo-phason term, which is a referee question we cannot answer with a citation and must answer with a
derivation and a statement of its size.

Timeline, fourteen weeks. Weeks one to two, reading, including the 2026 paper, and Step 1. Weeks three to
four, Step 2 and the branch count. Weeks five to six, Step 3. Weeks seven to eight, Step 4
and the two-cutoff map. Weeks nine to ten, Step 5. Weeks eleven to twelve, Step 6 and the sweeps. Weeks
thirteen to fourteen, figures and writing.

\begin{thebibliography}{99}
\small
\bibitem{Lubensky1985} T.~C.~Lubensky, S.~Ramaswamy, J.~Toner, \emph{Phys.\ Rev.\ B} \textbf{32} (1985) 7444, \doi{10.1103/PhysRevB.32.7444}.
\bibitem{Levine1985} D.~Levine, T.~C.~Lubensky, S.~Ostlund, S.~Ramaswamy, P.~J.~Steinhardt, J.~Toner, \emph{Phys.\ Rev.\ Lett.} \textbf{54} (1985) 1520, \doi{10.1103/PhysRevLett.54.1520}.
\bibitem{Agiasofitou2014} E.~Agiasofitou, M.~Lazar, \emph{Int.\ J.\ Solids Struct.} \textbf{51} (2014) 923, \doi{10.1016/j.ijsolstr.2013.11.016}.
\bibitem{Baggioli2020} M.~Baggioli, M.~Landry, \emph{SciPost Phys.} \textbf{9} (2020) 062, \doi{10.21468/SciPostPhys.9.5.062}.
\bibitem{Zhang2021} B.~Zhang, J.~G.~Yu, X.~M.~Zhang, L.~Elmaimouni, \emph{Acta Mech.} \textbf{232} (2021) 135, \doi{10.1007/s00707-020-02811-7}.
\bibitem{Wang2025} X.~Wang, J.~Yu, B.~Zhang, L.~Elmaimouni, P.~Ming, \emph{Wave Motion} \textbf{134} (2025) 103479, \doi{10.1016/j.wavemoti.2024.103479}.
\bibitem{Wang2022} X.~X.~Wang, J.~G.~Yu, B.~Zhang, L.~Elmaimouni, X.~M.~Zhang, X.~H.~Wang, \emph{Acta Mech.} \textbf{233} (2022) 3021, \doi{10.1007/s00707-022-03274-8}.
\bibitem{Wang2025b} X.~Wang, J.~Yu, B.~Zhang, L.~Elmaimouni, L.~Li, \emph{Compos.\ Struct.} \textbf{370} (2025) 119430, \doi{10.1016/j.compstruct.2025.119430}.
\bibitem{Zhang2026} B.~Zhang, J.~Song, L.~Li, H.~Han, J.~Yu, X.~Yang, \emph{Acta Mech.\ Solida Sin.} (2026), \doi{10.1007/s10338-026-00755-4}.
\bibitem{Chellappan2015} V.~Chellappan, S.~Gopalakrishnan, V.~Mani, \emph{J.\ Appl.\ Phys.} \textbf{117} (2015), \doi{10.1063/1.4907212}.
\bibitem{Lim2015} C.~W.~Lim, G.~Zhang, J.~N.~Reddy, \emph{J.\ Mech.\ Phys.\ Solids} \textbf{78} (2015) 298, \doi{10.1016/j.jmps.2015.02.001}.
\bibitem{Barretta2018} R.~Barretta, F.~Marotti de Sciarra, \emph{Int.\ J.\ Eng.\ Sci.} \textbf{130} (2018) 187, \doi{10.1016/j.ijengsci.2018.05.009}.
\end{thebibliography}
\end{document}
